% Generated draft: requires mathematical and source review.
\documentclass[11pt]{article}
\usepackage[margin=25mm]{geometry}
\usepackage{fontspec}
\setmainfont{Latin Modern Roman}
\usepackage{amsmath,amssymb,mathtools,unicode-math}
\setmathfont{Latin Modern Math}
\usepackage{enumitem,needspace,xurl,hyperref}
\setlength{\parindent}{0pt}
\setlength{\parskip}{5pt}
\setlength{\emergencystretch}{3em}
\newcommand{\N}{\mathbb{N}}
\newcommand{\Z}{\mathbb{Z}}
\newcommand{\Q}{\mathbb{Q}}
\newcommand{\R}{\mathbb{R}}
\newcommand{\C}{\mathbb{C}}
\newcommand{\F}{\mathbb{F}}
\newcommand{\A}{\mathbb{A}}
\newcommand{\PP}{\mathbb{P}}
\newcommand{\E}{\mathbb{E}}
\newcommand{\eps}{\varepsilon}
\newcommand{\dd}{\,\mathrm{d}}
\begin{document}
\section*{UnsolvedMath ID: 30005623}
\textbf{Source problem code:} catalogue entry
\subsection{Definitions and mathematical statement}
Let \(T\subset\R^2\) be the open interior of a nondegenerate triangle whose three interior angles are strictly less than \(\pi/2\), and let \(V(T)\) be its set of three vertices. List the Neumann eigenvalues of \(-\Delta\), counting multiplicity, as
\[
0=\lambda_1(T)<\lambda_2(T)\leq\lambda_3(T)\leq\cdots.
\]
Under the usual real-valued convention, a second Neumann eigenfunction is any nonzero \(u\in H^1(T;\R)\) satisfying
\[
\int_T\nabla u\cdot\nabla\varphi\,\dd x
=\lambda_2(T)\int_Tu\varphi\,\dd x
\qquad\text{for every }\varphi\in H^1(T;\R).
\]
Thus the boundary condition is the homogeneous Neumann condition \(\partial_\nu u=0\) along each open edge. No simplicity assumption on \(\lambda_2(T)\) is imposed; the quantifier includes every nonzero element of its eigenspace.

Interpreting a nonvertex critical point as a point of the closed triangle other than a vertex at which the full gradient vanishes, set
\[
C_T(u)=\{p\in\overline T\setminus V(T):\nabla u(p)=0\}.
\]
At an open-edge point, the gradient is evaluated using the smooth local extension supplied by reflection across that straight edge. Consequently, this convention counts both interior critical points and critical points on open edges.

The reconstructed assertion supplied by the registry is
\[
\text{for every acute planar triangle }T\text{ and every second Neumann eigenfunction }u,
\qquad \#C_T(u)\leq 1.
\]

\emph{Unresolved source qualifications.} The supplied statement does not explicitly specify the scalar field or whether open-edge points count as critical points. The formulation above makes the conventional real-valued, closed-triangle interpretation explicit; these conventions require confirmation against the intended source. Counting only interior points would give a different formulation. Moreover, the supplied original extract omits the preceding assertion to which its words ``the above statement'' refer and then changes to an unrelated graph problem. The registry labels its reconstructed triangle statement as corrected and verified, but the supplied extract alone does not establish that reconstruction. This is a statement-only draft requiring human mathematical review.
\subsection{Short English statement}
The supplied reconstructed claim says that every second Neumann eigenfunction on every acute planar triangle has at most one critical point away from the vertices. The displayed formulation uses real-valued eigenfunctions and counts points in the interior and on open edges. Those conventions and the reconstruction need review because the original extract omits the relevant preceding claim and includes an unrelated graph discussion.
\subsection{Sources}
\begin{itemize}
\item Local supplied UnsolvedMath registry record, permanent ID \(30005623\), source problem code \texttt{catalogue entry}, titled ``Critical Points of Second Neumann Eigenfunctions on Acute Triangles,'' in the Oberwolfach Reports Open Problems collection. The record gives proposed year 2023 and no proposer for the triangle problem. Its supplied status is \texttt{open}; its statement status is \texttt{corrected\_verified}, with the supplied assertion that the reconstruction was checked against the official report. These are registry assertions, not independent verification here.
\item Supplied report citation: \emph{Geometric Spectral Theory (2024)}, Oberwolfach Reports, pp.\ 2017--2106, DOI \texttt{10.4171/owr/2023/36}. Source URL and literature-source link labelled ``original OWR report'': \url{https://doi.org/10.4171/owr/2023/36}. The supplied title year 2024 is retained alongside the DOI containing 2023.
\item Supplied literature source labelled ``Judge--Mondal, Euclidean triangles'': \url{https://doi.org/10.4007/annals.2020.191.1.3}. No further bibliographic details or precise theorem from this reference are supplied.
\item Supplied literature source labelled ``Burdzy--Werner, hot-spots counterexample'': \url{https://doi.org/10.2307/121027}. No further bibliographic details or precise theorem from this reference are supplied.
\item Supplied literature assessment dated \texttt{2026-08-22}: the omitted preceding hot-spots claim cannot be identified or settled from the known counterexamples and positive classes cited in the curation; the target is therefore described as unclear and classified as open. This is supplied dated curation, not independently verified literature assessment. The background expressly requests independent verification of time-sensitive status and cited claims.
\item The supplied original extract says that an unspecified conjecture is false for non-simply connected planar domains, citing local reference [5], and reports an alternative proof of the hot spots conjecture for all acute triangles in 2021, citing local reference [11]. The bibliography entries for [5] and [11] and the identity of the speakers denoted by ``we'' are not supplied; neither reference number is identified here with a listed DOI.
\end{itemize}
\end{document}
